\documentclass{article}
\usepackage{amsmath,amssymb}
\begin{document}

@@DOCUMENT_HEADER@@

@@SECTION_1_HEADER@@

\textbf{Câu 1.} Trong hệ đơn vị SI, đơn vị của gia tốc là

@@TAB@@\textbf{A.} $\text{m}\cdot\text{s}^2$.@@TAB@@\textbf{B.} $\text{m}\cdot\text{s}$.@@TAB@@\textbf{C.} $\text{m}\cdot\text{s}^{-1}$.@@TAB@@\textbf{D.} $\text{m}\cdot\text{s}^{-2}$.

\textbf{Câu 2.} Đồng hồ đo thời gian hiện số MC964 khi đặt ở chế độ MODE A $+$ B dùng để

@@TAB@@\textbf{A.} đo thời gian vật chắn cổng quang điện nối với ổ B.

@@TAB@@\textbf{B.} đo tổng hai khoảng thời gian vật chắn cổng quang điện nối với ổ A và cổng quang điện nối với ổ B.

@@TAB@@\textbf{C.} đo thời gian vật chuyển động từ cổng quang điện nối với ổ A tới cổng quang điện nối với ổ B.

@@TAB@@\textbf{D.} đo thời gian vật chắn cổng quang điện nối với ổ A.

\textbf{Câu 3.} Chọn phát biểu đúng về vận tốc.

@@TAB@@\textbf{A.} Vận tốc là đại lượng vô hướng, luôn có giá trị không âm.

@@TAB@@\textbf{B.} Vận tốc là đại lượng vectơ, có hướng ngược với hướng của độ dịch chuyển.

@@TAB@@\textbf{C.} Vận tốc là đại lượng vô hướng, có thể nhận giá trị âm hoặc dương.

@@TAB@@\textbf{D.} Vận tốc là đại lượng vectơ, có hướng trùng với hướng của độ dịch chuyển.

\textbf{Câu 4.} Đối tượng nghiên cứu của Vật lí là

@@TAB@@\textbf{A.} các dạng vận động của vật chất và năng lượng.

@@TAB@@\textbf{B.} sự phát triển sinh học của vật chất hữu cơ.

@@TAB@@\textbf{C.} sự hình thành và phát triển của các lí thuyết vật lí.

@@TAB@@\textbf{D.} cuộc đời và sự nghiệp của các nhà vật lí.

\textbf{Câu 5.} Chuyển động nào sau đây được coi là rơi tự do?

@@TAB@@\textbf{A.} Vận động viên đang nhảy dù từ máy bay xuống.

@@TAB@@\textbf{B.} Viên sỏi nhỏ được thả rơi từ trên cao trong không khí.

@@TAB@@\textbf{C.} Chiếc lá khô đang rơi từ trên cành xuống.

@@TAB@@\textbf{D.} Sợi chỉ mảnh được thả rơi trong không khí.

\textbf{Câu 6.} Phát biểu nào sau đây về chuyển động thẳng biến đổi đều là \textbf{sai}?

@@TAB@@\textbf{A.} Trong chuyển động nhanh dần đều, vectơ gia tốc cùng chiều với chiều chuyển động.

@@TAB@@\textbf{B.} Trong chuyển động thẳng biến đổi đều, gia tốc không đổi theo thời gian.

@@TAB@@\textbf{C.} Trong chuyển động chậm dần đều, gia tốc luôn có giá trị âm.

@@TAB@@\textbf{D.} Trong chuyển động chậm dần đều, vectơ gia tốc ngược chiều với chiều chuyển động.

\textbf{Câu 7.} Gọi $\overline{A}$ là giá trị trung bình, $\Delta A'$ là sai số dụng cụ, $\overline{\Delta A}$ là sai số ngẫu nhiên và $\Delta A = \overline{\Delta A} + \Delta A'$ là sai số tuyệt đối của phép đo đại lượng $A$. Sai số tỉ đối của phép đo là

@@TAB@@\textbf{A.} $\delta A = \dfrac{\overline{\Delta A}}{\overline{A}} \cdot 100\%$.@@TAB@@\textbf{B.} $\delta A = \dfrac{\overline{A}}{\overline{\Delta A}} \cdot 100\%$.@@TAB@@\textbf{C.} $\delta A = \dfrac{\Delta A'}{\overline{A}} \cdot 100\%$.@@TAB@@\textbf{D.} $\delta A = \dfrac{\Delta A}{\overline{A}} \cdot 100\%$.

\textbf{Câu 8.} Đồ thị độ dịch chuyển -- thời gian $(d - t)$ của chuyển động thẳng đều có dạng là

@@TAB@@\textbf{A.} đường thẳng.@@TAB@@\textbf{B.} đường tròn.@@TAB@@\textbf{C.} đường hypebol.@@TAB@@\textbf{D.} đường parabol.

\textbf{Câu 9.} Một vật chuyển động thẳng biến đổi đều với vận tốc ban đầu $v_0$, sau thời gian $t$ vật có vận tốc $v$. Gia tốc của vật là

@@TAB@@\textbf{A.} $a = \dfrac{v_0 - v}{t}$.@@TAB@@\textbf{B.} $a = \dfrac{v_0 + v}{t}$.@@TAB@@\textbf{C.} $a = \dfrac{v^2 - v_0^2}{t}$.@@TAB@@\textbf{D.} $a = \dfrac{v - v_0}{t}$.

\textbf{Câu 10.} Một chất điểm đi được quãng đường $s$ trong khoảng thời gian $\Delta t$. Tốc độ trung bình của chất điểm là

@@TAB@@\textbf{A.} $v_{\text{tb}} = s \cdot \Delta t$.@@TAB@@\textbf{B.} $v_{\text{tb}} = s - \Delta t$.@@TAB@@\textbf{C.} $v_{\text{tb}} = \dfrac{s}{\Delta t}$.@@TAB@@\textbf{D.} $v_{\text{tb}} = \dfrac{\Delta t}{s}$.

\textbf{Câu 11.} Độ dịch chuyển của một vật cho biết

@@TAB@@\textbf{A.} độ dài và hướng của sự thay đổi vị trí của vật.

@@TAB@@\textbf{B.} vị trí và thời điểm chuyển động của vật.

@@TAB@@\textbf{C.} mức độ nhanh hay chậm của chuyển động.

@@TAB@@\textbf{D.} tổng độ dài quãng đường mà vật đã đi được.

\textbf{Câu 12.} Khi phòng thực hành xảy ra sự cố cháy, hành động nào sau đây là \textbf{sai}?

@@TAB@@\textbf{A.} Ngắt ngay toàn bộ nguồn điện của phòng thực hành.

@@TAB@@\textbf{B.} Di chuyển nhanh các chai lọ hóa chất, vật liệu dễ cháy ra khu vực an toàn.

@@TAB@@\textbf{C.} Không dùng bình khí $\text{CO}_2$ để dập đám cháy trên quần áo người hoặc đám cháy kim loại kiềm.

@@TAB@@\textbf{D.} Dùng nước dội trực tiếp để dập đám cháy ở nơi có thiết bị điện đang cắm nguồn.

@@SECTION_2_HEADER@@

\textbf{Câu 1.} Một bạn học sinh đạp xe từ cổng trường THPT Chuyên Quốc Học Huế đi thẳng theo đường Lê Lợi $1\text{ km}$, sau đó rẽ phải vào đường Hùng Vương đi thẳng $500\text{ m}$ để đến quán chè Hẻm như sơ đồ hình bên. Coi hai đoạn đường vuông góc với nhau.

@@CENTER_IMAGE_fig_cau13@@

@@TAB1@@\textbf{a)} Quãng đường đi được và độ lớn độ dịch chuyển của bạn học sinh là khác nhau.

@@TAB1@@\textbf{b)} Độ dịch chuyển của bạn học sinh là một đại lượng vô hướng.

@@TAB1@@\textbf{c)} Quãng đường bạn học sinh đã đi được là $1{,}5\text{ km}$.

@@TAB1@@\textbf{d)} Độ lớn độ dịch chuyển của bạn học sinh bằng $1180\text{ m}$.

\textbf{Câu 2.} Một vật chuyển động thẳng dọc theo trục $Ox$ có phương trình độ dịch chuyển $d = 4t - t^2$ ($d$ tính bằng m, $t$ tính bằng s, $t \ge 0$).

@@TAB1@@\textbf{a)} Vật chuyển động thẳng nhanh dần đều trong suốt quá trình chuyển động.

@@TAB1@@\textbf{b)} Ngay sau thời điểm $t = 0$, vật chuyển động theo chiều dương của trục $Ox$.

@@TAB1@@\textbf{c)} Gia tốc của vật bằng $1\text{ m/s}^2$.

@@TAB1@@\textbf{d)} Vận tốc của vật tại thời điểm $t = 1\text{ s}$ bằng $2\text{ m/s}$.

@@SECTION_3_HEADER@@

\textbf{Câu 1.} Tại giải vô địch điền kinh thế giới ở Berlin (Đức), vận động viên Usain Bolt lập kỉ lục thế giới chạy $100\text{ m}$ với thời gian $9{,}58\text{ s}$. Tốc độ trung bình của Usain Bolt trên đường chạy đó là bao nhiêu $\text{m/s}$? (Làm tròn kết quả đến chữ số hàng phần mười).

@@ANSWER_BOX@@

\textbf{Câu 2.} Hình bên là đồ thị độ dịch chuyển -- thời gian $(d - t)$ của một ô tô chuyển động thẳng. Vận tốc của ô tô bằng bao nhiêu $\text{m/s}$?

@@CENTER_IMAGE_fig_cau16@@

@@ANSWER_BOX@@

\textbf{Câu 3.} Một người đang lái xe máy với tốc độ $28{,}8\text{ km/h}$ thì nhìn thấy một cái hố trước mặt nên hãm phanh, xe chuyển động chậm dần đều và dừng lại sau $5\text{ s}$. Chọn chiều dương là chiều chuyển động. Gia tốc của xe là bao nhiêu $\text{m/s}^2$?

@@ANSWER_BOX@@

\textbf{Câu 4.} Một vật được thả rơi tự do từ độ cao $80\text{ m}$ so với mặt đất. Lấy $g = 10\text{ m/s}^2$. Thời gian rơi của vật là bao nhiêu giây?

@@ANSWER_BOX@@

@@SECTION_4_HEADER@@

\textbf{Câu 1.} Lúc 7 giờ, một ô tô chuyển động thẳng đều với tốc độ $40\text{ km/h}$ đi từ $A$ qua $B$ đến $C$. Biết $AB = 30\text{ km}$, $BC = 70\text{ km}$ ($A$, $B$, $C$ thẳng hàng theo thứ tự đó). Chọn gốc tọa độ tại $A$, chiều dương từ $A$ đến $B$, gốc thời gian lúc 7 giờ. Vẽ đồ thị độ dịch chuyển -- thời gian $(d - t)$ của ô tô.

\textbf{Câu 2.} Một vật chuyển động thẳng biến đổi đều dọc theo trục $Ox$ có đồ thị vận tốc -- thời gian $(v - t)$ như hình bên. Tính:

@@TAB1@@\textbf{a)} Gia tốc của vật.

@@TAB1@@\textbf{b)} Quãng đường vật đi được trong $10\text{ s}$ đầu tiên.

@@CENTER_IMAGE_fig_cau20@@

\textbf{Câu 3.} Thả rơi tự do một vật từ độ cao $100\text{ m}$ so với mặt đất. Lấy $g = 10\text{ m/s}^2$. Tính:

@@TAB1@@\textbf{a)} Quãng đường vật rơi được sau $3\text{ s}$ kể từ lúc thả.

@@TAB1@@\textbf{b)} Độ cao của vật so với mặt đất tại thời điểm đó.



@@HET@@

\end{document}
